\section*{Chapter \ref{ch:iv1:timed-test-9} --- Timed Test 9}
\markboth{\protect\omTocMark{}{Answers: Timed Test 9}}{\protect\omTocMark{}{Answers: Timed Test 9}}

\begin{solution}{iq:iv1:timed-test-9:1}
The expected number is $\frac83$.

Let $G_1$ and $G_2$ be the rounds of each coin's first head. We want $\mathbb{E}\max(G_1, G_2)$. Using $\max + \min = G_1 + G_2$\index{maximum and minimum}, the sum has mean 4.

The minimum is the first round in which either coin shows heads. A round has this with chance $\frac34$, so the minimum has mean $\frac43$. So the maximum has mean $4 - \frac43 = \frac83$.

A second route is by states: from the start, $E = 1 + \frac14E + \frac12 \cdot 2$, which gives $\frac34E = 2$, so $E = \frac83$. A simulation agrees.

\iqwrong{Answering 4 by adding the two means, or 2 by saying that they run side by side. The true value lies between the two.}
\end{solution}

\begin{solution}{iq:iv1:timed-test-9:2}
The probability is $\frac17$.

There are $\binom83 = 56$ triples. Any equilateral triangle must have three equal sides, and the distances between vertices are the edge, the face diagonal and the space diagonal.

Only the face diagonal works\index{counting by symmetry}. Each vertex has 3 others at a face-diagonal distance, and these three are also mutually at that distance. So any two of them form an equilateral triangle with the vertex, and each vertex lies in $\binom32 = 3$ triangles. Each triangle has 3 vertices, so there are $\frac{8 \times 3}{3} = 8$ triangles. Then $\frac{8}{56} = \frac17$.

An enumeration of the 56 triples confirms 8. No triple of edges closes up, as a cube has no triangle of its edges.

\iqwrong{Counting only 2 triangles, or the triangles of one tetrahedron, of which there are 4. The cube holds two such tetrahedra.}
\end{solution}

\begin{solution}{iq:iv1:timed-test-9:3}
The posterior mean is 1 and the variance is $\frac12$. With noise variance 4, they are $0.4$ and $0.8$.

Precisions add\index{conjugate normal prior}. The prior has precision 1, and the measurement has precision 1. So the posterior precision is 2, the variance is $\frac12$, and the mean is the precision-weighted average: $\frac{1 \cdot 0 + 1 \cdot 2}{2} = 1$.

With noise variance 4, the measurement has precision $\frac14$. Then the posterior precision is $\frac54$, the variance is $0.8$, and the mean is $\frac{\frac14 \cdot 2}{\frac54} = 0.4$.

A numerical integration of the posterior agrees. Limit checks: with no noise, the mean is 2, and with infinite noise, it stays 0.

\iqwrong{Giving a posterior variance above the prior's. A new measurement can only reduce it, in this normal case.}
\end{solution}

\begin{solution}{iq:iv1:timed-test-9:4}
The probability is $\frac{16}{49}$.

Since $10^6 = 2^6 5^6$, the divisors are $2^a5^b$ with $a$ and $b$ each from 0 to 6, which gives $7 \times 7 = 49$\index{divisor counting}.

A multiple of $10^3$ needs $a \ge 3$ and $b \ge 3$. That is 4 values for each exponent, so 16 divisors. Hence $\frac{16}{49}$.

A direct list of the divisors confirms the counts. The chance of $a \ge 3$ alone is $\frac47$, and the two exponents are independent, so the answer is its square. Note the pairing: a divisor $d$ matches $10^6/d$, and $10^3$ sits in the middle of the list.

\iqwrong{Using the chance $\frac{1}{1000}$, as for a random whole number. Divisors are not spread evenly, so it is the exponents that must be counted.}
\end{solution}

\begin{solution}{iq:iv1:timed-test-9:5}
The probability is about $0.24$, and the conditional mean is 1.

$X - Y$ is normal with mean 0 and variance 2\index{sum of normals}. So $P(X - Y > 1) = P(Z > \frac{1}{\sqrt2}) = P(Z > 0.71) = 0.24$.

Given $X + Y = 2$, symmetry gives $\mathbb{E}[X \mid X + Y = 2] = \mathbb{E}[Y \mid X + Y = 2]$, and the two add up to 2. So each is 1.

The regression route agrees: the slope of $X$ on $X + Y$ is $\frac12$, and the mean of $X + Y$ is 0, so the value is $\frac12 \times 2$. A simulation gives both.

Limit check: for a target of 0 in place of 1, the chance is $\frac12$, as it should be.

\iqwrong{Using a variance of 1 for $X - Y$. Variances add for a difference too, so the scale is $\sqrt2$.}
\end{solution}

\begin{solution}{iq:iv1:timed-test-9:6}
The smaller has mean $\frac12$, the larger has mean $\frac32$, and the probability is $\frac{1}{2e} \approx 0.184$.

The minimum is exponential with rate 2\index{memoryless property}, so its mean is $\frac12$. After it, the other variable has a fresh exponential life, by the memoryless property. So the gap has mean 1, and the maximum has mean $\frac32$.

The sum checks: $\frac12 + \frac32 = 2 = \mathbb{E}(X + Y)$.

$X > Y + 1$ needs two things: $X$ outlasts $Y$, with chance $\frac12$, and then it lasts at least 1 more, with chance $e^{-1}$. A simulation returns $0.184$.

Limit check: replacing 1 by 0 gives $\frac12$, and a large gap makes the chance vanish.

\iqwrong{Treating the gap as depending on the size of $Y$. The memoryless property makes it independent of how long the first has lasted.}
\end{solution}

\begin{solution}{iq:iv1:timed-test-9:7}
$P(N \text{ even}) = \frac{1 + e^{-2}}{2} \approx 0.568$, and the chance of no sale is $e^{-3} \approx 0.050$.

Add the series for $e^{\lambda}$ and $e^{-\lambda}$\index{parity trick}: the odd terms cancel. So $\sum_{k \text{ even}}\frac{\lambda^k}{k!} = \cosh\lambda$, and $P(\text{even}) = e^{-\lambda}\cosh\lambda = \frac{1 + e^{-2\lambda}}{2}$. At $\lambda = 1$, this is $0.568$. As $\lambda$ grows, it tends to $\frac12$, and at $\lambda = 0$ it is 1.

Thinning a Poisson process\index{Poisson thinning} with a chance of $\frac12$ gives a Poisson process at rate 3. So no sale has chance $e^{-3}$. Summing over the number of customers gives the same result: $\sum e^{-6}\frac{6^k}{k!}2^{-k} = e^{-6}e^{3}$. A simulation agrees.

\iqwrong{Using $e^{-6}$, which is the chance of no customers at all. A customer who does not buy is still no sale.}
\end{solution}

\begin{solution}{iq:iv1:timed-test-9:8}
The probability is $\frac{1800}{7776} = \frac{25}{108} \approx 0.231$.

Choose the two values that pair up: $\binom62 = 15$. Choose the single value: 4. Arrange the five dice\index{multinomial coefficient}: $\frac{5!}{2!\,2!\,1!} = 30$. That gives $15 \times 4 \times 30 = 1800$, out of $6^5 = 7776$.

An enumeration over all the 7776 outcomes confirms the count.

Note that choosing the two pair values is unordered. Doing it in order would count each outcome twice.

\iqwrong{Using $6 \times 5 \times 4 \times 30 = 3600$, which treats (3, 3, 5, 5, 2) and (5, 5, 3, 3, 2) as different results.}

\iqfollow{Is this more or less likely than one pair?}{One pair has $6 \times \binom53 \times \frac{5!}{2!} = 3600$ outcomes, or $0.463$, so it is twice as likely.}
\end{solution}

\begin{solution}{iq:iv1:timed-test-9:9}
The weights are 2 to 1, or $\frac23$ and $\frac13$, and the Sharpe ratio is about $0.707$.

The best portfolio is proportional to $\Sigma^{-1}\mu$\index{maximum Sharpe portfolio}. With no correlation, the weights are proportional to $\mu_i/\sigma_i^2$: $\frac{0.05}{0.01} = 5$ and $\frac{0.10}{0.04} = 2.5$. That is 2 to 1.

Each asset alone has a Sharpe ratio of $\frac12$. For uncorrelated assets, the squared Sharpe ratios add, so the best is $\sqrt{\frac14 + \frac14} = 0.707$.

Direct check: the portfolio has a return of $\frac23 \cdot 5 + \frac13 \cdot 10 = 6.67\%$, and a variance of $\frac49 \cdot 100 + \frac19 \cdot 400 = 88.9$, so the volatility is $9.43\%$. The ratio is $\frac{6.67}{9.43} = 0.707$.

\iqwrong{Choosing equal weights, as the two Sharpe ratios match. Equal Sharpe calls for equal risk, not equal capital.}
\end{solution}

\begin{solution}{iq:iv1:timed-test-9:10}
The probability is about $0.0475$.

Choose the two ranks: $\binom{13}{2} = 78$. Choose the two suits in each\index{counting hands}: $\binom42^2 = 36$. Choose the fifth card from the 44 cards of other ranks. The count is $78 \times 36 \times 44 = 123\,552$. Divided by the number of hands, this is $0.0475$.

A simulation returns the same. Note the 44: only cards of the other 11 ranks are allowed, and each has 4 suits.

As a rough check, $78 \times 36 \times 44 \approx 80 \times 1600 = 128\,000$, and that over $2.6$ million is about one hand in 20.

\iqwrong{Using $13 \times 12$ for the ranks, which counts each hand twice. The two pairs have no order, so the value doubles to about $0.095$.}
\end{solution}

\begin{solution}{iq:iv1:timed-test-9:11}
The last passenger gets their own seat with probability $\frac12$, and the ninth is displaced with probability $\frac13$.

Whenever someone sits at random, three things can happen\index{symmetry}. They may choose the first passenger's seat, and all is settled. They may choose the last person's seat, and the last person is displaced. Or they may choose another seat, and the choice passes to someone else. The first and second are equally likely at each step, so the last passenger's chance is $\frac12$.

For the ninth, the same holds with three seats left to matter: seat 1, seat 9, and seat 10. When a random choice first falls on one of these, each is equally likely. Only seat 9 displaces the ninth passenger. So the chance is $\frac13$. In general, passenger $k$ of $n$ is displaced with chance $\frac{1}{n - k + 2}$.

A simulation confirms both. With $k = 10$, this is $\frac12$, as before.

\iqwrong{Computing the full tree of 10 passengers. The answer for the last passenger does not depend on $n$ at all, and that is the clue to look for a symmetry.}
\end{solution}

\begin{solution}{iq:iv1:timed-test-9:12}
$100!$ has 158 digits, and the limit is $\frac1e \approx 0.368$.

With $n = 100$, $n\ln n = 100 \times 4.605 = 460.5$\index{Stirling's approximation}. Subtracting 100 gives $360.5$. And $\frac12\ln(628) = 3.22$, so $\ln 100! \approx 363.7$. Dividing by $2.303$ gives about $157.9$. Its whole part is 157, so $100!$ has 158 digits. An exact product confirms it: $100! \approx 9.3 \times 10^{157}$.

The answer is close to the next power of ten, so a careful estimate matters here: the $\frac12\ln(2\pi n)$ term is essential.

For the limit, $\frac1n\ln n! \approx \ln n - 1$. So $(n!)^{1/n} \approx \frac{n}{e}$. This means that the geometric mean of $1, \dots, n$ is about $\frac{n}{e}$, against the arithmetic mean of about $\frac{n}{2}$.

\iqwrong{Dropping the small term and getting $\frac{360.5}{2.303} = 156.5$, so 157 digits. This is too low by one.}
\end{solution}
